% this file is called up by thesis.tex
% content in this file will be fed into the main document

%: ----------------------- introduction file header -----------------------
\chapter{Introduction}

% the code below specifies where the figures are stored
\ifpdf
    \graphicspath{{1_introduction/figures/PNG/}{1_introduction/figures/PDF/}{1_introduction/figures/}}
\else
    \graphicspath{{1_introduction/figures/EPS/}{1_introduction/figures/}}
\fi

% ----------------------------------------------------------------------
%: ----------------------- introduction content ----------------------- 
% ----------------------------------------------------------------------
\section{Background and Context}
Chronic conditions—including cardiovascular disease, diabetes, and frailty—are leading contributors to morbidity and mortality in ageing populations~\cite{bencivengaAtrialFibrillationElderly2020, bardonPredictorsIncidentMalnutrition2020}. These conditions are associated with declines in functional capacity (e.g., gait speed, grip strength~\cite{davisComparisonGaitSpeed2021}), increased healthcare utilisation~\cite{mooreCountingTimeLived2017}, and reduced quality of life~\cite{huangSexDifferencesAssociations2020}.  

Traditional risk prediction models rely primarily on clinical biomarkers such as blood pressure and lipid profiles. Evidence suggests that multidimensional predictors, including nutrition~\cite{murphyPlasmaLuteinZeaxanthin2023}, physical activity~\cite{lairdPhysicalActivityDose2023}, mental health~\cite{mcdowellAssociationsPhysicalActivity2018}, and social determinants~\cite{mccrorySocialDisadvantageSocial2016}, improve predictive performance.

The Irish Longitudinal Study on Ageing (\gls{TILDA}) provides a nationally representative dataset of adults aged 50+ in Ireland with repeated measures across multiple domains. The Harmonized \gls{TILDA} dataset aligns variables with the RAND Health and Retirement Study (\gls{HRS}) framework for cross-country comparability, while the raw waves contain additional study-specific variables. The IDS dataset supplements coverage for participants with disabilities.  

These data motivate the development of integrated models combining traditional (e.g., blood pressure, cholesterol, diabetes diagnosis) and non-traditional (e.g., nutrition biomarkers, physical activity, mental health, social determinants) risk factors. This thesis focuses on the non-traditional predictors to quantify their added predictive value.

% ----------------------------------------------------------------------
% Research Problem
\section{Research Problem and Motivation}
Existing \gls{ML} approaches for chronic disease prediction in older adults face several challenges:

\subsection{Limited Integration of Multidimensional Data}
Most models prioritise clinical variables but neglect modifiable lifestyle and nutrition factors~\cite{murphyPlasmaLuteinZeaxanthin2023, lairdPhysicalActivityDepression2023}.

\subsection{Underuse of Longitudinal Dynamics}
Few studies exploit repeated measures to capture trajectories of risk, such as frailty progression~\cite{moguilnerImportanceAgePrediction2021}.

\subsection{Lack of Explainability}
Many \gls{ML} models are not interpretable, limiting clinical utility~\cite{donoghueUsingConditionalInference2023}.

Addressing these gaps requires methods that integrate multidomain predictors, account for longitudinal dynamics, and balance predictive accuracy with interpretability.

% ----------------------------------------------------------------------
% Research Aims
\section{Research Aims and Objectives}
This thesis develops and validates \gls{ML} models that integrate nutrition, lifestyle, and mental health features with clinical data to improve prediction of chronic disease and ageing-related outcomes.

\begin{itemize}
  \item Conduct a literature review on \gls{ML} for chronic disease prediction in ageing populations, emphasising gaps in lifestyle and nutrition integration.
  \item Preprocess and harmonise \gls{TILDA} waves 1--5 and engineer multidomain features (nutrition, activity, mobility, mental health, clinical indicators).
  \item Develop and compare baseline statistical models (logistic regression) with advanced \gls{ML} methods (\gls{RF}, \gls{GBM}).
  \item Evaluate models using discrimination (\gls{AUC}), calibration, and error metrics (accuracy, precision, recall, F1-score); perform ablation analyses to quantify added value of lifestyle and nutrition.
  \item Identify actionable predictors (e.g., vitamin D status, gait speed reserve) and highlight modifiable risk factors for intervention.
  \item Document ethical, governance, and reproducibility practices for handling longitudinal pseudonymised health data.
\end{itemize}

% ----------------------------------------------------------------------
% Methodology Overview (concise, non-repetitive)
\section{Methodology Overview}
This section summarises the workflow; detailed methods are presented in Chapter 3.  

Data are derived from both raw \gls{TILDA} waves 1--5 and the Harmonized \gls{TILDA} dataset, aligned to the \gls{HRS} framework. IDS data complements coverage for participants with disabilities.  

Missing data are addressed with complete-case analysis and multiple imputation. Derived features capture multidimensional predictors:
\begin{itemize}
    \item Frailty indices (e.g., Fried’s phenotype, \gls{FI})
    \item Mobility measures (e.g., \gls{TUG}, gait speed, grip strength)
    \item Mental health scales (e.g., \gls{CESD}, \gls{GDS})
    \item Nutrition scores (e.g., dietary diversity indices, micronutrient biomarkers)
\end{itemize}

Baseline (logistic regression) and advanced \gls{ML} models (\gls{RF}, \gls{GBM}) are tuned via cross-validation. Performance is evaluated using discrimination (\gls{AUC}), calibration (plots, Brier score), and error metrics. Explainability is assessed with feature importance and \gls{SHAP} values.

% ----------------------------------------------------------------------
% Contributions
\section{Contributions of the Research}
\begin{itemize}
  \item \textbf{Empirical:} Quantifies the predictive value of nutrition and lifestyle features beyond clinical baselines.
  \item \textbf{Methodological:} Provides a reproducible pipeline for harmonising longitudinal \gls{TILDA} data and integrating multidomain features into \gls{ML} models.
  \item \textbf{Clinical:} Highlights modifiable predictors for targeted interventions and policy to promote healthy ageing.
\end{itemize}

% ----------------------------------------------------------------------
% Thesis Outline
\section{Thesis Outline}
\begin{itemize}
    \item \emph{Chapter Two:} Literature review on \gls{ML} prediction of chronic disease in ageing populations.
    \item \emph{Chapter Three:} \gls{TILDA} study design, data preprocessing, harmonisation, and feature engineering.
    \item \emph{Chapter Four:} Modelling framework, including baseline and advanced \gls{ML}, hyperparameter tuning, validation, and evaluation metrics.
    \item \emph{Chapter Five:} Results, model comparison, feature importance, and ablation analyses.
    \item \emph{Chapter Six:} Discussion of implications, limitations, ethical considerations, and future work.
\end{itemize}

\vspace{5 mm}

A series of documents have been included in the Appendix section of this dissertation. These are:

\begin{itemize}
\item \emph{Appendix A}: Variable codebook and wave harmonisation notes (including anonymisation actions, top-coding, and grouping rules relevant to the features used).
\item \emph{Appendix B}: Ethics.
\item \emph{Appendix C}: Modelling details, additional tables/figures, and reproducibility checklist (overview of data processing scripts).
\end{itemize} 

\vspace{5 mm}

A digital archive is attached to this dissertation containing:(Access request required)

\begin{itemize}
\item \emph{folder 1}: Preprocessing and feature-engineering scripts (documentation only; no raw \gls{TILDA} data).
\item \emph{folder 2}: Model training/evaluation code and configuration files to reproduce reported results.
\end{itemize}